% GENERATED FILE. Edit canonical lessons, then run npm run generate:books.
% canonical-source-hash: fnv1a64:ba87e0a52c322de9
% canonical-lessons: HI-C158-caura, HI-C158-gahra, HI-C158-uca, HI-C158-nica, HI-C158-khali

\chapter{Wide, Deep, High, Low, Empty}
\label{ch:hi-wide-deep-high-low-empty}
\hlchaptermodality{158}

\begin{chapteropening}
\textbf{By the end of this chapter:} \emph{I can say wide, deep, high, low and empty.}

Complete the last lesson of chapter 158: I can say wide, deep, high, low and empty.
\end{chapteropening}

\section[\texorpdfstring{\hi{चौड़ा} (cauṛā)}{चौड़ा (cauṛā)}]{\hi{चौड़ा} (cauṛā) — wide}

\label{lesson:HI-C158-caura}



\begin{quote}
Before the new one: say the Hindi for young, then the Hindi for old.
\end{quote}

\subsection*{You'll want to know: \hi{चौड़ा}}
\textbf{\hi{चौड़ा}} — \emph{cauṛā} — \textquotedblleft{}wide\textquotedblright{}.

\begin{cousinweb}
\textbf{\hi{चौड़ा}} (\emph{cauṛā}) is a tadbhava, a word inherited and reshaped by sound change. It comes from \emph{*ca-uḍa}, \textquotedblleft{}wide, open\textquotedblright{}, a form CDIAL reconstructs from its descendants (CDIAL 4534). CDIAL suggests, with a \textquotedblleft{}possibly\textquotedblright{}, that it was built on \emph{cau-}, \textquotedblleft{}four\textquotedblright{}.
\end{cousinweb}

\begin{grammarlens}[title={What you've built}]
One more word for this chapter.
\end{grammarlens}

\subsection*{Your turn}
\begin{itemize}
\raggedright
  \item \emph{Say it:} \emph{cauṛā}
  \item \emph{Say it:} \emph{cauṛā}, once more
  \item \emph{Say it:} say \emph{cauṛā}
  \item \emph{Recall:} say the Hindi for young, then the Hindi for old, then say \emph{cauṛā} again
\end{itemize}

\begin{tcolorbox}[breakable,colback=teal!4,colframe=teal!35!black,title={Before you move on}]
What does \hi{चौड़ा} mean? (\textquotedblleft{}Wide\textquotedblright{}.) Say it once more.
\end{tcolorbox}
\section[\texorpdfstring{\hi{गहरा} (gahrā)}{गहरा (gahrā)}]{\hi{गहरा} (gahrā) — deep, dark (of colour)}

\label{lesson:HI-C158-gahra}



\begin{quote}
Before the new one: say the Hindi for old, then the Hindi for wide.

Say the colours and their stories: \textbf{\hi{लाल}}, a red gem's name; \textbf{\hi{नीला}}, indigo's dark blue; \textbf{\hi{सफ़ेद}}, a Persian loan.
\end{quote}

\subsection*{You'll want to know: \hi{गहरा}}
\textbf{\hi{गहरा}} — \emph{gahrā} — \textquotedblleft{}deep, dark (of colour)\textquotedblright{}.

\begin{cousinweb}
\textbf{\hi{गहरा}} (\emph{gahrā}) is a tadbhava, a word inherited from Sanskrit and reshaped by sound change. It comes from \emph{gabhīra}, \textquotedblleft{}deep, solemn\textquotedblright{} (CDIAL 4024).
\end{cousinweb}

\begin{grammarlens}[title={What you've built}]
One more word for this chapter.
\end{grammarlens}

\subsection*{Your turn}
\begin{itemize}
\raggedright
  \item \emph{Say it:} \emph{gahrā}
  \item \emph{Say it:} \emph{gahrā}, once more
  \item \emph{Say it:} say \emph{gahrā}
  \item \emph{Recall:} say the Hindi for old, then the Hindi for wide, then say \emph{gahrā} again
\end{itemize}

\begin{tcolorbox}[breakable,colback=teal!4,colframe=teal!35!black,title={Before you move on}]
What does \hi{गहरा} mean? (\textquotedblleft{}Deep, dark (of colour)\textquotedblright{}.) Say it once more.
\end{tcolorbox}
\section[\texorpdfstring{\hi{ऊँचा} (ū̃cā)}{ऊँचा (ū̃cā)}]{\hi{ऊँचा} (ū̃cā) — high, tall}

\label{lesson:HI-C158-uca}



\begin{quote}
Before the new one: say the Hindi for wide, then the Hindi for deep.
\end{quote}

\subsection*{You'll want to know: \hi{ऊँचा}}
\textbf{\hi{ऊँचा}} — \emph{ū̃cā} — \textquotedblleft{}high, tall\textquotedblright{}.

\begin{cousinweb}
\textbf{\hi{ऊँचा}} (\emph{ūṁcā}) is a tadbhava, a word inherited from Sanskrit and reshaped by sound change. It comes from \emph{ucca}, \textquotedblleft{}tall, lofty\textquotedblright{} (CDIAL 1634). CDIAL calls the influence of \emph{udañc}, \textquotedblleft{}turned upwards\textquotedblright{}, on the nasal very doubtful.
\end{cousinweb}

\begin{grammarlens}[title={What you've built}]
One more word for this chapter.
\end{grammarlens}

\subsection*{Your turn}
\begin{itemize}
\raggedright
  \item \emph{Say it:} \emph{ū̃cā}
  \item \emph{Say it:} \emph{ū̃cā}, once more
  \item \emph{Say it:} say \emph{ū̃cā}
  \item \emph{Recall:} say the Hindi for wide, then the Hindi for deep, then say \emph{ū̃cā} again
\end{itemize}

\begin{tcolorbox}[breakable,colback=teal!4,colframe=teal!35!black,title={Before you move on}]
What does \hi{ऊँचा} mean? (\textquotedblleft{}High, tall\textquotedblright{}.) Say it once more.
\end{tcolorbox}
\section[\texorpdfstring{\hi{नीचा} (nīcā)}{नीचा (nīcā)}]{\hi{नीचा} (nīcā) — low}

\label{lesson:HI-C158-nica}



\begin{quote}
Before the new one: say the Hindi for deep, then the Hindi for high.
\end{quote}

\subsection*{You'll want to know: \hi{नीचा}}
\textbf{\hi{नीचा}} — \emph{nīcā} — \textquotedblleft{}low\textquotedblright{}.

\begin{cousinweb}
\textbf{\hi{नीचा}} (\emph{nīcā}) is a tadbhava, a word inherited from Sanskrit and reshaped by sound change. It comes from \emph{nīca}, \textquotedblleft{}low\textquotedblright{} (CDIAL 7540). CDIAL explains its \emph{-c-} as shaped after its opposite, \emph{ucca}, \textquotedblleft{}high\textquotedblright{}.
\end{cousinweb}

\begin{grammarlens}[title={What you've built}]
One more word for this chapter.
\end{grammarlens}

\subsection*{Your turn}
\begin{itemize}
\raggedright
  \item \emph{Say it:} \emph{nīcā}
  \item \emph{Say it:} \emph{nīcā}, once more
  \item \emph{Say it:} say \emph{nīcā}
  \item \emph{Recall:} say the Hindi for deep, then the Hindi for high, then say \emph{nīcā} again
\end{itemize}

\begin{tcolorbox}[breakable,colback=teal!4,colframe=teal!35!black,title={Before you move on}]
What does \hi{नीचा} mean? (\textquotedblleft{}Low\textquotedblright{}.) Say it once more.
\end{tcolorbox}
\section[\texorpdfstring{\hi{ख़ाली} (khālī)}{ख़ाली (khālī)}]{\hi{ख़ाली} (khālī) — empty}

\label{lesson:HI-C158-khali}



\begin{quote}
Before the new one: say the Hindi for high, then the Hindi for low.
\end{quote}

\subsection*{You'll want to know: \hi{ख़ाली}}
\textbf{\hi{ख़ाली}} — \emph{khālī} — \textquotedblleft{}empty\textquotedblright{}.

\begin{cousinweb}
\textbf{\hi{ख़ाली}} (\emph{khālī}) is a loan from Arabic \emph{khālī}, \textquotedblleft{}empty\textquotedblright{}, which reached Hindi through Persian. Its Arabic root is \emph{kh-l-w} (Lane).
\end{cousinweb}

\begin{grammarlens}[title={What you've built}]
One more word for this chapter.
\end{grammarlens}

\subsection*{Your turn}
\begin{itemize}
\raggedright
  \item \emph{Say it:} \emph{khālī}
  \item \emph{Say it:} \emph{khālī}, once more
  \item \emph{Say it:} say \emph{khālī}
  \item \emph{Recall:} say the Hindi for high, then the Hindi for low, then say \emph{khālī} again
\end{itemize}

\begin{tcolorbox}[breakable,colback=teal!4,colframe=teal!35!black,title={Before you move on}]
What does \hi{ख़ाली} mean? (\textquotedblleft{}Empty\textquotedblright{}.) Say it once more.
\end{tcolorbox}
